\documentclass[11pt,a4paper]{article}
\usepackage[T1]{fontenc}
\usepackage{lmodern,amsmath,amssymb,amsthm,mathtools,microtype}
\usepackage[margin=27mm]{geometry}
\usepackage[colorlinks=true,linkcolor=blue,citecolor=blue,urlcolor=blue]{hyperref}
\setlength{\emergencystretch}{2em}
\newtheorem{theorem}{Theorem}[section]
\newtheorem{lemma}[theorem]{Lemma}
\newtheorem{corollary}[theorem]{Corollary}
\newcommand{\Q}{\mathbb Q}
\title{Square-complement normalization does not preserve the pole predicate}
\author{}\date{16 November 2026}
\begin{document}\maketitle
\begin{abstract}
The rational substitution $\Phi(b)=1156b/(17+b)^2$ preserves input
square class and makes $17(17-\Phi(b))$ a rational square. We prove that
it does not preserve the normalized pole predicate. For the infinite
family $b_n=204^2/409^{2n-1}$, $n\ge2$, no parameter gives a pole point
over $\Q_{409}$, including all zero and projective boundary cases.
The transformed inputs are fully prepared and have explicit nonzero
rational witnesses at every finite nonzero rational conic parameter. The result
concerns the isolated pole subsystem; no full detector-source
counterexample or defining-polynomial lower bound is asserted.
\end{abstract}

\section{The normalized equations and the substitution}

For a nonzero rational input $b=dt^2$, with $d,t\ne0$, put
\begin{equation}\label{eq:data}
 D=17d,\qquad C=17(17-b),\qquad
 e_\tau=\frac{\tau^2-D}{\tau^2+3D},\qquad A_\tau=e_\tau C.
\end{equation}
The normalized first conic and second pole of \cite{UnitNorm,Returns} are
\begin{equation}\label{eq:pole}
 p^2=D(1-e_\tau)(1+3e_\tau),\qquad
 q^2=D(s^2-A_\tau)(s^2+3A_\tau).
\end{equation}
The first conic has the ordinate
\begin{equation}\label{eq:p}
                    p_\tau=-\frac{4D\tau}{\tau^2+3D}.
\end{equation}
The missing conic parameter $\tau=\infty$ gives $e=1$, $p=0$.
Indeed every first-conic solution with $e\ne1$ has
$\tau=p/(e-1)$ and is recovered by these formulas.

A prepared rational input has square reduced denominator and satisfies
\[
 0<b<17,\qquad
 b\in\Q_2^{\times2}\cap\Q_3^{\times2}\cap\Q_{17}^{\times2},
 \qquad v_2(b)\ge3,\quad v_3(b)\ge1,\quad v_{17}(b)\ge2.
\]
No prepared-input existence theorem is needed for the obstruction below.

For $b\ne0,-17$, define
\begin{equation}\label{eq:phi}
 \Phi(b)=\frac{1156b}{(17+b)^2}
        =b\left(\frac{34}{17+b}\right)^2.
\end{equation}
Direct expansion gives
\begin{equation}\label{eq:square}
 17(17-\Phi(b))=
       \left(\frac{17(17-b)}{17+b}\right)^2.
\end{equation}
Thus the rational square class of the input is unchanged.
Its denominator information need not be unchanged: for $p\nmid34$
and $v_p(b)<0$ one has
\begin{equation}\label{eq:reverse}
                         v_p(\Phi(b))=-v_p(b)>0.
\end{equation}

\begin{samepage}
\begin{theorem}\label{thm:obstruction}
For $n\ge2$ let
\[
 P_n=409^{2n-1},\qquad b_n=\frac{204^2}{P_n},\qquad B_n=\Phi(b_n).
\]
Then:
\begin{enumerate}
\item $0<b_n<1$, and $b_n$ satisfies all the distinguished local-square
      and valuation conditions above. Its reduced denominator is not
      a square: $v_{409}(b_n)=1-2n<0$ is odd.
\item For $b_n$, no local conic parameter in $\mathbf P^1(\Q_{409})$
      gives a point on the smooth projective second-pole curve. This
      includes $s=0$, $q=0$, and the points above $s=\infty$.
\item The transformed input $B_n$ is prepared. With the same
      square-class representative $d=409$, it has nonzero rational
      witnesses for both equations in~\eqref{eq:pole} at every
      rational $\tau\ne0$.
\end{enumerate}
In particular \eqref{eq:phi} is not an equivalence of the isolated
pole predicates.
\end{theorem}
\end{samepage}

\section{Uniform local nonexistence before normalization}

The prime $409$ satisfies
\[
 409\equiv1\pmod8,\qquad409\equiv1\pmod3,\qquad
 409\equiv1\pmod{17},\qquad204=2^2\cdot3\cdot17.
\]
It is a square in each of $\Q_2,\Q_3,\Q_{17}$, by the usual
local square criteria. Therefore $b_n$ is a square in all three fields,
with valuations $4,2,2$. Since $409^3>204^2$, it lies in $(0,1)$.
Its numerator is prime to $409$, proving part~1 of the theorem.
It is not prepared; hence its local nonexistence does not conflict
with the local theorem for prepared inputs in \cite{Local}.

Use the decomposition
\[
        b_n=409\left(\frac{204}{409^n}\right)^2,
        \qquad D=17\cdot409=6953.
\]
Write $v=v_{409}$. For every $\tau\in\Q_{409}$, the valuations
of $\tau^2$ and $D$ are different: one is even and the other is one.
With $v(0)=+\infty$, this includes $\tau=0$. Since $3$ is a unit,
\[
 v(\tau^2-D)=v(\tau^2+3D)=\min\{2v(\tau),1\}.
\]
Both expressions are nonzero. Thus $e_\tau$ is a nonzero unit.
Since $v(b_n)=1-2n<0$, it follows that
\begin{equation}\label{eq:A}
                 v(C)=v(A_\tau)=1-2n.
\end{equation}
For every $s\in\Q_{409}$, the even valuation of $s^2$ differs
from this odd integer, again with the convention at zero. Hence
\[
 v(s^2-A_\tau)=v(s^2+3A_\tau)
                  =\min\{2v(s),1-2n\},
\]
and both factors are nonzero. Their pole radicand has odd valuation
\begin{equation}\label{eq:odd}
 v\bigl(D(s^2-A_\tau)(s^2+3A_\tau)\bigr)
                  =1+2\min\{2v(s),1-2n\}.
\end{equation}
It is not a square. This excludes every affine solution, including
$s=0$ or $q=0$.

The chart $u=1/s$, $z=q/s^2$ at infinity has equation
\[
               z^2=D(1-A_\tau u^2)(1+3A_\tau u^2).
\]
At $u=0$ it would require $z^2=D$, impossible because $v(D)=1$.
The quartic has four distinct roots over an algebraic closure:
$DA_\tau\ne0$, and $A_\tau\ne-3A_\tau$ in characteristic zero.
Thus the two charts account for the smooth projective completion.
For the conic parameter $\tau=\infty$, set $e=1$; the same
valuation argument applies. This proves part~2.

The obstruction is independent of the rational representative $d$ of
the input square class. Any such representative has odd $409$-adic
valuation. Replace the leading one in~\eqref{eq:odd} by that odd
valuation; the argument, including both boundaries, is unchanged.

\section{Prepared inputs and explicit witnesses after normalization}

Direct cancellation in~\eqref{eq:phi} gives
\begin{equation}\label{eq:B}
                B_n=\frac{408^2P_n}{(P_n+2448)^2}.
\end{equation}
This fraction is reduced. Indeed $2448=6\cdot408$ and
$409\equiv1\pmod{408}$, so
\[
 \gcd(P_n+2448,408)=1,
 \qquad \gcd(P_n+2448,P_n)=1.
\]
The reduced denominator is therefore a square, and
$v_{409}(B_n)=2n-1>0$ is odd. Identity~\eqref{eq:square} gives
$0<B_n<17$, because $0<b_n<1$.
The quotient $B_n/b_n$ is a nonzero rational square, so the three
local square conditions are retained. Moreover
$P_n+2448\equiv1\pmod{408}$ and $408=2^3\cdot3\cdot17$, giving
\[
                    v_2(B_n)=6,\qquad v_3(B_n)=2,
                    \qquad v_{17}(B_n)=2.
\]
Thus $B_n$ is fully prepared. Its decomposition in the same square
class is
\[
 B_n=409\left(\frac{408\,409^{n-1}}{P_n+2448}\right)^2.
\]

Put
\[
             c_n=\frac{17(P_n-2448)}{P_n+2448}>0.
\]
Then $17(17-B_n)=c_n^2$. For any rational $\tau\ne0$, retain
$D=6953$, take $e=e_\tau$ and $p=p_\tau$ from
\eqref{eq:data}, \eqref{eq:p}, and set
\begin{equation}\label{eq:witness}
                         s=c_n,\qquad q=c_n^2p_\tau.
\end{equation}
The second-pole equation is the exact identity
\[
 D(s^2-e c_n^2)(s^2+3e c_n^2)
          =c_n^4D(1-e)(1+3e)=c_n^4p_\tau^2=q^2.
\]
All denominators are nonzero since $D>0$. In addition
\[
 1-e_\tau=\frac{4D}{\tau^2+3D},\qquad
 1+3e_\tau=\frac{4\tau^2}{\tau^2+3D}.
\]
Thus both factors and both ordinates are nonzero. At the single
uniform parameter $\tau=1$, the witness is
\[
 e=-\frac{1738}{5215},\qquad p=-\frac{6953}{5215},\qquad
 s=c_n,\qquad q=-\frac{6953}{5215}c_n^2.
\]
This proves part~3 and Theorem~\ref{thm:obstruction}.

The same substitution~\eqref{eq:witness} works whenever $D>0$ and
$C=c^2\ne0$, with $c_n$ replaced by $c$. Consequently a square
complement can supply automatic pole witnesses while the transformation
that produced it removes the forbidden negative odd denominator.

\newpage
\section{Subsystem scope and exact verification}

The family above has not been realized as actual detector values for a
nonintegral input in a full defining-polynomial source. The theorem
excludes treating \eqref{eq:phi} as an unconditional equivalence of pole
predicates; it does not exclude an independently proved construction
which recovers the original denominator information by other conditions.

The standard-library checker is
\begin{center}\small
\texttt{negative\_results/computation/}\\
\texttt{square\_complement\_normalization\_v43/check.py}.
\end{center}
It verifies the polynomial identities, primality and congruence
constants, coprimality formulas, exact family instances, and both witness
equations. It also checks the valuation alternatives on a finite audit
grid and on exact rational local examples including zero parameters.
The proof above supplies the universal quantifiers and both projective
boundary arguments; the finite checks do not replace those arguments.

\begin{thebibliography}{9}
\raggedright
\bibitem{UnitNorm}
\emph{A single norm determinant for the 25- and 13-variable constructions},
project source, version 28,
\texttt{papers/finite\_algebra/}\allowbreak
\texttt{unit\_norm\_20261114\_v28.tex}.
\bibitem{Returns}
\emph{Fixed pole fibers and independent detector indices},
project source, version 42,
\texttt{papers/pole\_parameters/}\allowbreak
\texttt{common\_parameter\_guard\_reduction\_20261116\_v42.tex}.
\bibitem{Local}
\emph{Local solubility at every shared pole parameter},
project source, version 41,
\texttt{papers/pole\_parameters/}\allowbreak
\texttt{shared\_parameter\_local\_solubility\_20261115\_v41.tex}.
\end{thebibliography}
\end{document}
